\section{Related Work}
\label{sec:related_work}

\para{Activation steering.}
The linear representation hypothesis~\citep{park2023linear,marks2023geometry}
posits that high-level concepts correspond to approximately linear directions
in transformer activation space. \actadd~\citep{turner2023actadd} exploits this
by computing a steering vector from the activation difference between a positive
and negative prompt pair and adding it at a single layer during the forward pass.
\caa~\citep{panickssery2024caa} extends this to use mean-difference vectors from
hundreds of contrastive pairs, applied at all token positions after the
instruction. The design choice of continuous injection implicitly acknowledges
that single-application steering would decay, though neither work measures this
directly. Representation Engineering~\citep{zou2023repe} provides a broader
framework for reading and writing model representations, treating steering as
manipulating a learned representation space.

\para{Evidence for steering decay.}
Several papers provide indirect evidence that steering effects diminish during
generation. \citet{li2026svf} demonstrate that static steering vectors become
misaligned during autoregressive generation because hidden states drift through
activation space, causing reliability to degrade in long-form generation. They
propose Steering Vector Fields (\svf), which re-compute the steering direction
from the current hidden state every $K$ steps. \citet{nguyen2025pid} formalize
steering as a control theory problem and prove that proportional-only control
(equivalent to standard activation addition) has inherent steady-state error.
Their PID controller adds integral and derivative terms to counteract this error.
\citet{bigelow2025belief} model steering through a Bayesian lens, where the
steering vector shifts a concept prior that can be overridden by evidence
accumulated during generation. Unlike these works, which propose solutions
to decay, we directly measure the decay dynamics themselves.

\para{Mechanisms of decay.}
Several mechanisms have been proposed for how perturbations are absorbed.
\citet{xu2026why} argue that model activations lie on a low-dimensional manifold,
and subsequent layers project perturbed activations back onto this manifold.
\citet{dang2026selective} identify norm distortion as a specific mechanism:
adding a steering vector violates activation norms expected by LayerNorm,
causing distribution shift that compounds across layers. They propose
norm-preserving rotation as an alternative. The Hidden Life of
Tokens~\citep{hiddentokens2025} shows that information injected into hidden
states diminishes as more tokens are generated, suggesting that perturbations
are naturally attenuated by the autoregressive process.

\para{Behavioral persistence and error propagation.}
\citet{bachmann2024pitfalls} establish that small perturbations in hidden states
can compound through autoregressive generation due to the teacher-forcing gap.
This suggests competing effects: the steering perturbation may decay in hidden
states, but behavioral changes caused by early steered tokens may persist as
the model conditions on its own outputs.
Our experimental design explicitly separates these two effects by comparing
autoregressive generation (where both operate) with teacher-forced conditions
(where only representational decay is measured).

\para{Counteracting decay.}
Beyond continuous injection~\citep{panickssery2024caa},
methods include context-dependent refreshing~\citep{li2026svf},
PID control~\citep{nguyen2025pid},
conditional steering~\citep{lee2024cast},
semantics-adaptive intervention~\citep{wang2024sadi},
norm-preserving rotation~\citep{dang2026selective},
SAE-targeted steering~\citep{chalnev2024saesteer},
and KV cache steering~\citep{kvcache2025}.
Our work complements these by characterizing the problem they address:
we show that without continuous reinforcement, steering influence is lost
within a single step, motivating the need for these approaches.

\para{Evaluation of steering.}
\citet{pres2024reliable} demonstrate that multiple-choice evaluations can
overstate steering success and that long-form generation is a stricter test.
\citet{tan2024steerability} and \citet{bas2025limits} show that steerability
varies across concepts and follows an inverted-U curve with steering strength.
These findings motivate our choice to evaluate across multiple behaviors and to
measure both short-term (hidden state) and long-term (output distribution)
effects.
